%% kapa: exam=1 round=36 year=2025 period=1 form=A subject=real_estate subject_ko=부동산학원론 seq=20 number=100
%% source: Q-Net 감정평가사 기출문제 (cst00309) artlSeq=5231525 file=raw/1008/5231525/2025년 제36회 감정평가사 1차시험 1교시 문제지 원본.pdf
%% conversion: kordoc 4.24.51 -> tools/gen1.py (auto); review=proofread-sonnet; flags=none
\begin{question}{100}
대출상환에 관한 설명으로 옳지 않은 것은? (단, 주어진 조건에 한함)
\begin{choices}
  \item 대출조건이 동일할 경우, 대출채권의 듀레이션(duration)은 원리금균등분할상환, 원금균등분할상환, 점증상환, 만기일시상환의 순으로 짧다.
  \item 원리금균등분할상환의 경우, 매월 원리금을 균등하게 상환하기 때문에 원리금에서 원금과 이자가 차지하는 비중은 상환시기에 따라 다르다.
  \item 점증상환에서는 초기에 원리금의 납입액이 이자지급액에 미치지 못할 수 있는데, 이 경우 미상환 이자가 원금에 가산되어 부(-)의 상환이 일어날 수 있다.
  \item 만기일시상환은 대출기간 동안 매월 이자만 상환하다가 만기에 일시로 원금을 상환하는 방식이며, 대출만기 시 원금의 일부를 상환하게 한 뒤 대출만기를 연장해주기도 한다.
  \item 원금균등분할상환의 경우, 매월 상환하는 원리금 상환 부담은 대출 초기에는 많지만 상환금액은 점차 감소한다.
\end{choices}
\end{question}
\answer{1}
